% TOP_PROBLEM: {"id":30007577,"problem_number":"LOCAL-30007577","title":"Trust in complex auction rules","category_id":21,"category":{"id":21,"name":"economics","display_name":"Economics","description":"Open economic theory statements, published research agendas, and proposed scoped specifications; question provenance and historical priority are qualified per record.","slug":"economics","order_index":21,"created_at":"2026-10-08T00:00:00Z"},"status":"research_question","external_url":null,"legacy_ids":[],"published":false,"statement":"Identify how explanations and credible rule verification affect participation, bidder losses, and auction welfare, while separating measured comprehension and trust channels under stated assumptions.","economics":{"original_id":66,"field":"Mechanism design and market design","kind":"empirical","formalization_basis":"adapted_research_question","source_status":"research_agenda","priority_status":"earliest_found","priority_note":"This is the earliest source in which the relevant agenda was verified during this audit; absolute priority has not been established.","earliest_verified_reference":{"authors":["Ignacio Palacios-Huerta","David C. Parkes","Richard Steinberg"],"title":"Combinatorial Auctions in Practice","year":2024,"url":"https://eprints.lse.ac.uk/124108/1/Combinatorial_Auctions_in_Practice.pdf","doi":"10.1257/jel.20221679","locator":"§4.3 Simplicity and Trust on Both Sides of the Market, manuscript pp.23–25 (PDF pp.24–26)","evidence_summary":"The authors identify future research on practical simplicity, trust, and credibility, alongside established mechanism-design and implementation results."},"current_reference":{"authors":["Ignacio Palacios-Huerta","David C. Parkes","Richard Steinberg"],"title":"Combinatorial Auctions in Practice","year":2024,"url":"https://eprints.lse.ac.uk/124108/1/Combinatorial_Auctions_in_Practice.pdf","doi":"10.1257/jel.20221679","locator":"§4.3 Simplicity and Trust on Both Sides of the Market, manuscript pp.23–25 (PDF pp.24–26)","evidence_summary":"The authors identify future research on practical simplicity, trust, and credibility, alongside established mechanism-design and implementation results."},"formalization_note":"The randomization, causal outcomes, and mediation qualifications are our empirical adaptation of §4.3; theoretical credibility itself has existing results."},"statusQualification":"Published question or research agenda verified at the cited locator. The mathematical specification is adapted for this catalogue and includes additional assumptions; source publication does not establish first-ever priority or certify this exact model remains unsolved. The randomization, causal outcomes, and mediation qualifications are our empirical adaptation of §4.3; theoretical credibility itself has existing results.","statusReviewedAt":"2026-10-08"}
% FIRST_POSTED: 2026-10-08
% ENTRY_KIND: statement_only
% !TeX program = lualatex
\documentclass[11pt]{article}
\usepackage{fontspec}
\usepackage[a4paper,margin=25mm]{geometry}
\usepackage{amsmath,amssymb,mathtools}
\usepackage{xurl}
\usepackage[hidelinks]{hyperref}
\newcommand{\sref}[2]{\hyperlink{source:#1:#2}{[S#2]}}
\setlength{\emergencystretch}{3em}
\begin{document}
% problemId:30007577
\section*{Trust in complex auction rules}\label{30007577}
\subsection{Definitions and mathematical statement}
Fix a combinatorial auction with goods \(G\), allocation rule \(x\), payment rule \(p\), and bidder values drawn from a stated population distribution. Randomize groups across interfaces \(Z\) that differ in explanations, rule verification, or commitment evidence while implementing identical economic allocation and payment rules. Let \(Y_i(Z)\) denote bidder participation, \(L_i(Z)\) the loss relative to the bidder's best response under the observed auction environment, and \(W(Z)\) total allocative surplus. Elicit perceived probabilities \(b_i(Z)\) that the auctioneer follows the announced rule, using an incentivized method with specified measurement error. Estimate
\[
 E[Y_i(1)-Y_i(0)],\quad E[L_i(1)-L_i(0)],\quad E[W(1)-W(0)].
\]
Determine how the treatment works through perceived credibility, comprehension, and participation costs. Any mediation decomposition must state its additional identification assumptions; randomization of an explanation alone does not identify the separate causal effect of trust. Include auctioneer-side deviations in a separate prespecified design study, with deviations classified by detectability and the auctioneer's revenue incentives. Evaluate effects on experienced as well as inexperienced bidders and test persistence across auctions. Established credible-mechanism theory and cryptographic verification remain background. The open agenda is whether implementable communication and commitment devices produce portable participation and efficiency gains in complex auctions, with experimentally credible attribution of the mechanism.

\par\textbf{Formulation and scope.} The randomization, causal outcomes, and mediation qualifications are our empirical adaptation of §4.3; theoretical credibility itself has existing results.

\par\textbf{Open-question provenance.} An explicit question or research agenda was verified at the source locator. This does not by itself certify that every later version remains unresolved \sref{30007577}{1}.

\par\textbf{Historical priority.} This is the earliest source in which the relevant agenda was verified during this audit; absolute priority has not been established.

\subsection{Short English statement}
Identify how explanations and credible rule verification affect participation, bidder losses, and auction welfare, while separating measured comprehension and trust channels under stated assumptions.

\subsection{Sources}
\begin{itemize}\raggedright
\item[\textnormal{[S1]}]\hypertarget{source:30007577:1}{} Ignacio Palacios-Huerta, David C. Parkes, Richard Steinberg. 2024. Combinatorial Auctions in Practice. §4.3 Simplicity and Trust on Both Sides of the Market, manuscript pp.23–25 (PDF pp.24–26).
\url{https://eprints.lse.ac.uk/124108/1/Combinatorial_Auctions_in_Practice.pdf}.
DOI: 10.1257/jel.20221679.
{\small\textit{Earliest verified question/agenda reference; historical priority qualified below. The authors identify future research on practical simplicity, trust, and credibility, alongside established mechanism-design and implementation results.}}
\end{itemize}
\end{document}
